% GENERATED FILE. Edit canonical lessons, then run npm run generate:books.
% canonical-source-hash: fnv1a64:5acfa7ae25000ae0
% canonical-lessons: ZH-C236-chuxian, ZH-W236-xian-appear, ZH-C236-faxian, ZH-C236-haipa, ZH-C236-xinghao, ZH-C236-chadianr

\chapter{Something Odd}
\label{ch:zh-something-odd}
\hlchaptermodality{236}

\begin{chapteropening}
\textbf{By the end of this chapter:} \emph{I can say something turned up, someone found out, how scared I was, what luckily happened and what nearly did.}

Complete the last lesson of chapter 236: I can say something turned up, someone found out, how scared I was, what luckily happened and what nearly did.
\end{chapteropening}

\section[\texorpdfstring{\zh{出现} (chūxiàn)}{出现 (chūxiàn)}]{\zh{出现} (chūxiàn) — to appear; to turn up}

\label{lesson:ZH-C236-chuxian}



\begin{quote}
Before the new one: say the Mandarin for to have no choice but, then the Mandarin for no wonder.
\end{quote}

\subsection*{You'll want to know: \zh{出现}}
\textbf{\zh{出现}} — \emph{chūxiàn} — \textquotedblleft{}to appear; to turn up\textquotedblright{}.

\begin{quote}
\textbf{\zh{草地里出现了一个人}} \emph{cǎodì lǐ chūxiànle yí ge rén} — Someone appeared on the lawn.
\end{quote}

The new character, in pieces: \zh{现} — eight strokes in two pieces: on the left two bars, a vertical and a stroke rising to the right, then \zh{见} on the right. The next lesson puts it in your hand.

\begin{cousinweb}
\textbf{\zh{出现}} (\emph{chūxiàn}) is built from \zh{出} 'to come out' and \zh{现} 'to show', both inherited from Middle Chinese, \emph{tsyhwit} and \emph{henH} (Guangyun, in Baxter's transcription), which glosses the second 'to reveal', noting the character as an everyday form.
\end{cousinweb}

\begin{grammarlens}[title={What you've built}]
One more word for talking about the past.
\end{grammarlens}

\subsection*{Your turn}
\begin{itemize}
\raggedright
  \item \emph{Say it:} \emph{chūxiàn}
  \item \emph{Say it:} \emph{chūxiàn}, once more
  \item \emph{Say it:} \emph{cǎodì lǐ chūxiànle yí ge rén} — Someone appeared on the lawn.
  \item \emph{Recall:} say the Mandarin for to have no choice but, then the Mandarin for no wonder, then say \emph{chūxiàn} again
\end{itemize}

\begin{tcolorbox}[breakable,colback=teal!4,colframe=teal!35!black,title={Before you move on}]
What does \zh{出现} mean? (\textquotedblleft{}To appear; to turn up\textquotedblright{}.) Say it once more.
\end{tcolorbox}
\section[\texorpdfstring{\zh{现} (xiàn)}{现 (xiàn)}]{\zh{现} — a narrow left piece, then \zh{见}}

\label{lesson:ZH-W236-xian-appear}



\begin{quote}
The word first, then the character inside it.

\begin{itemize}
\raggedright
  \item \emph{Recall:} say \emph{chūxiàn} — \textbf{R1}, one lesson back, before it decays
\end{itemize}

\begin{quote}
\zh{见} \zh{在}
\end{quote}

\emph{Read it:} each one aloud, with nothing said first — \emph{jiàn}, \emph{zài}
\end{quote}

\subsection*{Script — a narrow left piece, then \zh{见}}
\hlblockfigure{\detokenize{figures/ZH-W236-xian-appear-filmstrip.pdf}}{How \zh{现} is written, stroke by stroke}

\begin{quote}
\zh{现}
\end{quote}

\textbf{Two pieces, left before right.}

\textbf{The left piece, four strokes,} narrow: a bar, a second bar, a vertical down through both, then a short stroke rising to the right at the foot.

\textbf{The right piece, \zh{见}, four strokes,} as in \zh{再见} and \zh{看见}: a vertical, a bar that runs right and turns down, a long stroke falling to the left, then a stroke that bends right along the foot and flicks up.

4 + 4 = eight.

\subsection*{Your turn}
\begin{itemize}
\raggedright
  \item \emph{Write it:} \zh{现}, all eight, with the model in front of you
  \item \emph{Cover:} write it again from memory, and count the pieces, not the strokes
\end{itemize}

\begin{tcolorbox}[breakable,colback=teal!4,colframe=teal!35!black,title={Before you move on}]
How many strokes in \zh{现}? (\textbf{Eight}.) What is the right piece? (\textbf{\zh{见}}.) Which way does the last stroke of the left piece go? (\textbf{It rises to the right}.)

Source: \href{https://raw.githubusercontent.com/chanind/hanzi-writer-data/68d10a4b21150cae5e1ebbd223eed289cf32d90c/data/\%E7\%8E\%B0.json}{Hanzi Writer Data \zh{现}, pinned stroke paths and medians}.
\end{tcolorbox}
\section[fāxiàn]{fāxiàn — to discover; to find out}

\label{lesson:ZH-C236-faxian}



\begin{quote}
Before the new one: say the Mandarin for no wonder, then the Mandarin for to turn up.
\end{quote}

\subsection*{You'll want to know: fāxiàn}
\emph{fāxiàn} — \textquotedblleft{}to discover; to find out\textquotedblright{}.

\begin{quote}
\emph{huí jiā yǐhòu wǒ fāxiàn qiánbāo diū le} — When I got home, I found my wallet was gone.
\end{quote}

Its characters come later; for now it is a word for the ear.

\begin{cousinweb}
\emph{fāxiàn} is built from \emph{fā} 'to send out' and \zh{现} 'to show', both inherited from Middle Chinese, \emph{pjot} and \emph{henH} (Guangyun, in Baxter's transcription), which glosses the first 'to rise' and 'to make clear'.
\end{cousinweb}

\begin{grammarlens}[title={What you've built}]
One more word for talking about the past.
\end{grammarlens}

\subsection*{Your turn}
\begin{itemize}
\raggedright
  \item \emph{Say it:} \emph{fāxiàn}
  \item \emph{Say it:} \emph{fāxiàn}, once more
  \item \emph{Say it:} \emph{huí jiā yǐhòu wǒ fāxiàn qiánbāo diū le} — When I got home, I found my wallet was gone.
  \item \emph{Recall:} say the Mandarin for no wonder, then the Mandarin for to turn up, then say \emph{fāxiàn} again
\end{itemize}

\begin{tcolorbox}[breakable,colback=teal!4,colframe=teal!35!black,title={Before you move on}]
What does fāxiàn mean? (\textquotedblleft{}To discover; to find out\textquotedblright{}.) Say it once more.
\end{tcolorbox}
\section[hàipà]{hàipà — to be scared}

\label{lesson:ZH-C236-haipa}



\begin{quote}
Before the new one: say the Mandarin for to turn up, then the Mandarin for to find out.
\end{quote}

\subsection*{You'll want to know: hàipà}
\emph{hàipà} — \textquotedblleft{}to be scared\textquotedblright{}.

\begin{quote}
\emph{wǎnshang yí ge rén zài jiā wǒ hěn hàipà} — Home alone at night, I was really scared.
\end{quote}

Its characters come later; for now it is a word for the ear.

\begin{cousinweb}
\emph{hàipà} is built from \emph{hài} 'harm' and \emph{pà} 'to fear', both inherited from Middle Chinese, \emph{hajH} and \emph{phaeH} (Guangyun, in Baxter's transcription), which glosses the first 'to injure'.
\end{cousinweb}

\begin{grammarlens}[title={What you've built}]
One more word for talking about the past.
\end{grammarlens}

\subsection*{Your turn}
\begin{itemize}
\raggedright
  \item \emph{Say it:} \emph{hàipà}
  \item \emph{Say it:} \emph{hàipà}, once more
  \item \emph{Say it:} \emph{wǎnshang yí ge rén zài jiā wǒ hěn hàipà} — Home alone at night, I was really scared.
  \item \emph{Recall:} say the Mandarin for to turn up, then the Mandarin for to find out, then say \emph{hàipà} again
\end{itemize}

\begin{tcolorbox}[breakable,colback=teal!4,colframe=teal!35!black,title={Before you move on}]
What does hàipà mean? (\textquotedblleft{}To be scared\textquotedblright{}.) Say it once more.
\end{tcolorbox}
\section[xìnghǎo]{xìnghǎo — luckily; fortunately}

\label{lesson:ZH-C236-xinghao}



\begin{quote}
Before the new one: say the Mandarin for to find out, then the Mandarin for to be scared.
\end{quote}

\subsection*{You'll want to know: xìnghǎo}
\emph{xìnghǎo} — \textquotedblleft{}luckily; fortunately\textquotedblright{}.

\begin{quote}
\emph{xìnghǎo nǐ lái le} — Thank goodness you came.
\end{quote}

Its characters come later; for now it is a word for the ear.

\begin{cousinweb}
\emph{xìnghǎo} is built from \emph{xìng} 'fortunate' and \zh{好} 'good', both inherited from Middle Chinese, \emph{heangX} and \emph{xawX} (Guangyun, in Baxter's transcription), which explains the first, after an older book, as getting the good and escaping the bad.
\end{cousinweb}

\begin{grammarlens}[title={What you've built}]
One more word for talking about the past.
\end{grammarlens}

\subsection*{Your turn}
\begin{itemize}
\raggedright
  \item \emph{Say it:} \emph{xìnghǎo}
  \item \emph{Say it:} \emph{xìnghǎo}, once more
  \item \emph{Say it:} \emph{xìnghǎo nǐ lái le} — Thank goodness you came.
  \item \emph{Recall:} say the Mandarin for to find out, then the Mandarin for to be scared, then say \emph{xìnghǎo} again
\end{itemize}

\begin{tcolorbox}[breakable,colback=teal!4,colframe=teal!35!black,title={Before you move on}]
What does xìnghǎo mean? (\textquotedblleft{}Luckily; fortunately\textquotedblright{}.) Say it once more.
\end{tcolorbox}
\section[chàdiǎnr]{chàdiǎnr — nearly; almost}

\label{lesson:ZH-C236-chadianr}



\begin{quote}
Before the new one: say the Mandarin for to be scared, then the Mandarin for luckily.
\end{quote}

\subsection*{You'll want to know: chàdiǎnr}
\emph{chàdiǎnr} — \textquotedblleft{}nearly; almost\textquotedblright{}.

\begin{quote}
\emph{wǒ chàdiǎnr chídào le} — I was nearly late.
\end{quote}

Its characters come later; for now it is a word for the ear.

\begin{cousinweb}
\emph{chàdiǎnr} is built from \emph{chà} 'short of' and \zh{点儿} 'a little', inherited from Middle Chinese \emph{tsrhae}, \emph{temX} and \emph{nye} (Guangyun, in Baxter's transcription), which glosses the first 'to choose' and 'to go wrong'; the falling tone of chà is not the regular outcome of that form.
\end{cousinweb}

\begin{grammarlens}[title={What you've built}]
One more word for talking about the past.
\end{grammarlens}

\subsection*{Your turn}
\begin{itemize}
\raggedright
  \item \emph{Say it:} \emph{chàdiǎnr}
  \item \emph{Say it:} \emph{chàdiǎnr}, once more
  \item \emph{Say it:} \emph{wǒ chàdiǎnr chídào le} — I was nearly late.
  \item \emph{Recall:} say the Mandarin for to be scared, then the Mandarin for luckily, then say \emph{chàdiǎnr} again
\end{itemize}

\begin{tcolorbox}[breakable,colback=teal!4,colframe=teal!35!black,title={Before you move on}]
What does chàdiǎnr mean? (\textquotedblleft{}Nearly; almost\textquotedblright{}.) Say it once more.
\end{tcolorbox}
